\vspace{-3mm}
\subsection{Equivariant Graph Attention}
\label{subsec:equivariant_graph_attention}
\vspace{-2mm}

Self-attention~\citep{transformer, gat, se3_transformer, transformer_in_vision, graphormer, gatv2} transforms features sent from one spatial location to another with input-dependent weights.
We use the notion from Transformers~\citep{transformer} and message passing networks~\citep{neural_message_passing_quantum_chemistry} and define message $m_{ij}$ sent from node $j$ to node $i$ as follows:
\begin{equation}
m_{ij} = a_{ij} \times v_{ij}
\end{equation}
%where attention weights $a_{ij}$ depend on features on node $i$ and its neighbors $\mathcal{N}(i)$ and values $v_{ij}$ are features sent from node $j$ to node $i$ transformed with input-independent weights.
where attention weights $a_{ij}$ depend on features on node $i$ and its neighbors $\mathcal{N}(i)$ and values $v_{ij}$ are transformed with input-independent weights.
In Transformers and Graph Attention Networks (GAT)~\citep{gat, gatv2}, $v_{ij}$ depends only on node $j$.
In message passing networks~\citep{neural_message_passing_quantum_chemistry}, $v_{ij}$ depends on features on nodes $i$ and $j$ with constant $a_{ij}$.
The proposed equivariant graph attention adopts tensor products to incorporate content and geometric information and uses multi-layer perceptron attention for $a_{ij}$ and non-linear message passing for $v_{ij}$ as illustrated in Fig.~\ref{fig:equiformer}(b).

\vspace{-2mm}
\paragraph{Incorporating Content and Geometric Information.}
Given features $x_i$ and $x_j$ on target node $i$ and source node $j$, we combine the two features with two linear layers to obtain initial message $x_{ij} = \text{Linear}_{dst} (x_i) + \text{Linear}_{src}(x_j)$.
$x_{ij}$ is passed to a DTP layer and a linear layer to consider geometric information like relative position contained in different type-$L$ vectors in irreps features: 
\begin{equation}
x_{ij}' = x_{ij} \otimes_{w(||\vec{r}_{ij}||)}^{DTP} \text{SH}(\vec{r}_{ij}) 
\quad\text{and}\quad
f_{ij} = \text{Linear}(x_{ij}')
\label{eq:pre_attn}
\end{equation}
where $x_{ij}'$ is the tensor product of $x_{ij}$ and spherical harmonics embeddings (SH) of relative position $\vec{r}_{ij}$, with weights parametrized by $||\vec{r}_{ij}||$.
%\revision{Note that the one-to-one dependence of channels in DTP significantly reduces the number of weights, and this enables generating weights } 
%To mix information in different channels, we use another linear layer following DTP and this results in message $f_{ij}$.
%The linear layer following DTP mixes information in difference channels and results in $f_{ij}$.
$f_{ij}$ considers semantic and geometric features on source and target nodes in a linear manner and is used to derive attention weights and non-linear messages.
%\todo{emphasize more on geometric information (background)}

\vspace{-2.25mm}
\paragraph{Multi-Layer Perceptron Attention.}
Attention weights $a_{ij}$ capture how each node interacts with neighboring nodes. 
$a_{ij}$ are invariant to geometric transformation, and thus, we only use type-$0$ vectors (scalars) of message $f_{ij}$ denoted as $f_{ij}^{(0)}$ for attention.
Note that $f_{ij}^{(0)}$ encodes directional information, as they are generated by tensor products of type-$L$ vectors with $L \geq 0$.
Inspired by GATv2~\citep{gatv2}, we adopts multi-layer perceptron attention (MLPA) instead of dot product attention (DPA) used in Transformers~\citep{transformer}.
In contrast to dot product, MLPs are universal approximators~\citep{mlp_universal_approximator, approximation_mlp, approximation_sigmoid} and can theoretically capture any attention patterns. 
% from theoretical perspectives
%%%Similar to GAT~\cite{gat, gatv2},
Given $f_{ij}^{(0)}$, we uses one leaky ReLU layer and one linear layer for $a_{ij}$:
\vspace{-1mm}
\begin{equation}
z_{ij} = a^\top \text{LeakyReLU}(f_{ij}^{(0)}) 
\quad\text{and}\quad
a_{ij} = \text{softmax}_{j}(z_{ij}) = \frac{\text{exp}(z_{ij})}{\sum_{k \in \mathcal{N}(i)} \text{exp}(z_{ik})}
\end{equation}
where $a$ is a learnable vectors of the same dimension as $f_{ij}^{(0)}$ and $z_{ij}$ is a single scalar.
%The output of attention is the sum of value $v_{ij}$ multipled by corresponding $a_{ij}$:
The output of attention is the sum of value $v_{ij}$ multipled by corresponding $a_{ij}$ over all neighboring nodes $j \in \mathcal{N}(i)$,
%\begin{equation}
%y_{i} = \sum_{j \in \mathcal{N}(i)} a_{ij} \times v_{ij}
%\end{equation}
where $v_{ij}$ can be obtained by linear or non-linear transformations of $f_{ij}$ as discussed below.

\vspace{-2.25mm}
\paragraph{Non-Linear Message Passing.}
Values $v_{ij}$ are features sent from one node to another, transformed with input-independent weights.
We first split $f_{ij}$ into $f_{ij}^{(L)}$ and $f_{ij}^{(0)}$, where the former consists of type-$L$ vectors with $0 \leq L \leq L_{max}$ and the latter consists of scalars only.
Then, we perform non-linear transformation to $f_{ij}^{(L)}$ to obtain non-linear message:
\begin{equation}
\mu_{ij} = \text{Gate}(f_{ij}^{(L)}) 
\quad\text{and}\quad
v_{ij} = \text{Linear}( [ \mu_{ij} \otimes_{w}^{DTP} \text{SH}(\vec{r}_{ij}) ] ) 
\end{equation}
%We obtain $\mu_{ij}$ by applying gate activation to $f_{ij}^{(L)}$.
We apply gate activation to $f_{ij}^{(L)}$ to obtain $\mu_{ij}$.
We use one DTP and a linear layer to enable interaction between non-linear type-$L$ vectors, which is similar to how we transform $x_{ij}$ into $f_{ij}$.
Weights $w$ here are input-independent.
We can also use $f_{ij}^{(L)}$ directly as $v_{ij}$ for linear messages.


\vspace{-2.25mm}
\paragraph{Multi-Head Attention.}
Following Transformers~\citep{transformer}, we can perform $h$ parallel equivariant graph attention functions given $f_{ij}$.
The $h$ different outputs are concatenated and projected with a linear layer, resulting in the final output $y_i$ as illustrated in Fig.~\ref{fig:equiformer}(b).
Note that parallelizing attention functions and concatenating can be implemented with ``Reshape''.